\documentclass[12pt]{article}
\usepackage[utf8]{inputenc}
\usepackage[T1]{fontenc}
\usepackage{hyperref}
\title{Fanbak Language Resources for K-12 Curriculum Development}
\author{Compiled by Claude Sonnet 5.5 (claude-sonnet-5-5) for Andr\'as Kornai}
\date{}
\begin{document}
\maketitle

\section{Introduction}
Fanbak (Fanbyak, also associated with Orkon) is a small Oceanic language of Ambrym Island, Vanuatu. Estimates of speakers range from about 30 active speakers to roughly 130; speakers are multilingual in North Ambrym and Bislama, and younger people in Port Vila tend to shift to Bislama. The main resources found are documentation-project materials.

\section{Primary Educational Resources}
The Newcastle news item below reports vital language resources created for Vanuatu communities; the search results did not show whether Fanbak school materials are included.
\begin{itemize}
\item Newcastle University news: \url{https://www.newcastle.edu.au/news/2017/05/two-year-project-culminates-in-the-creation-of-vital-language-resources-for-vanuatu-communities}
\end{itemize}

\section{Linguistic and Reference Materials}
\begin{itemize}
\item Fanbak language (Wikipedia): \url{https://en.wikipedia.org/wiki/Fanbak_language}
\item A Documentation of Orkon/Fanbak, a language of Vanuatu: \url{https://cultureincrisis.org/projects/a-documentation-of-orkon-fanbak-a-language-of-vanuatu}
\end{itemize}

\section{Religious and Community Resources}
None found.

\section{Digital and Audio Resources}
\begin{itemize}
\item The languages of northern Ambrym, Vanuatu: a guide to the deposited materials in ELAR: \url{https://scholarspace.manoa.hawaii.edu/handle/10125/24849}
\item The languages of Northern Ambrym, Vanuatu (ELDP archive): \url{https://eldp.access.preservica.com/dk01310387/}
\end{itemize}

\section{Curriculum Development Context}
The language is endangered and not known to be taught. The archived documentation materials could support community-led development of teaching materials; access conditions were not checked.
\end{document}